%% notes-tex/modeling/scaling-laws/sections/1-forms.tex

\section{The two forms\secstatus{todo}}
\label{sec:forms}

A \term{scaling law} is a fitted relation between a model's held-out loss and one
or more of the three quantities that a training run spends: the number of trainable
parameters $N$, the amount of training data $D$, and the training compute $C$. The
claim behind the name is that, over the range a family of models was trained on,
the loss falls as a power of each quantity, so that a line fit on log-log axes to
the runs one can afford predicts the loss of a run one cannot. The literature
carries two functional forms, and the difference between them is a floor.

\notesec{Kaplan: one axis at a time}
Kaplan et al.\ (arXiv:2001.08361)\external{arXiv:2001.08361, not in the mirror}
fit each axis separately, with the other axes made large enough not to limit:
\begin{equation}
  L(N) = \left(\frac{N_c}{N}\right)^{\alpha_N}, \qquad
  L(D) = \left(\frac{D_c}{D}\right)^{\alpha_D}.
  \label{eq:kaplan}
\end{equation}
Here $L$ is the \term{test loss}, the average loss on a fixed held-out evaluation
set, cross-entropy per token for a language model. $N_c$ and $D_c$ are fitted
constants in units of parameters and data; they set where each curve sits and have
no meaning of their own. $\alpha_N$ and $\alpha_D$ are the \term{scaling exponents},
the slopes on log-log axes, and the larger one names the axis along which a decade
of spending buys more. For language models Kaplan reports $\alpha_N \approx 0.076$
and $\alpha_D \approx 0.095$. Compute enters through the transformer rule
\begin{equation}
  C \approx 6 N D \ \text{FLOPs},
  \label{eq:sixnd}
\end{equation}
about two floating-point operations per parameter per training token for the
forward pass and four for the backward pass. $N$ in that rule and in
Equation~\ref{eq:kaplan} counts \emph{non-embedding} parameters; including the
embedding matrix bends the curve at small $N$, because the embedding scales with the
vocabulary rather than with the capacity the loss responds to.

\notesec{Chinchilla: a joint form with a floor}
Hoffmann et al.\ (arXiv:2203.15556)\external{arXiv:2203.15556, not in the mirror}
fit both axes at once and let the loss bottom out:
\begin{equation}
  L(N, D) = E + \frac{A}{N^{\alpha}} + \frac{B}{D^{\beta}}.
  \label{eq:chinchilla}
\end{equation}
$E$ is the \term{irreducible loss}, the value approached with infinite parameters
and infinite data; it comes from noise or inherent unpredictability in the data,
the entropy of text for a language model. $A$ and $B$ are the coefficients of the
model-size and data terms, and $\alpha$, $\beta$ are the exponents of the joint fit.
Their Approach~3 fit gives $E = 1.69$, $A = 406.4$, $B = 410.7$, $\alpha = 0.34$,
$\beta = 0.28$ on their training runs. With a floor present, a pure power law is
wrong at the large end: on log-log axes $L$ bends toward $E$ while $L - E$ stays
straight, which is the diagnostic panel~c of Figure~\ref{fig:forms} draws.

\notesec{Compute-optimal allocation}
Minimizing Equation~\ref{eq:chinchilla} subject to the budget $C = 6ND$ gives the
size and data a fixed budget should buy. Writing $D = C/(6N)$ and setting the
derivative in $N$ to zero,
\begin{equation}
  N^*(C) = G \left(\frac{C}{6}\right)^{a}, \qquad
  D^*(C) = \frac{C}{6 N^*(C)}, \qquad
  G = \left(\frac{\alpha A}{\beta B}\right)^{\frac{1}{\alpha + \beta}}, \qquad
  a = \frac{\beta}{\alpha + \beta}, \quad b = \frac{\alpha}{\alpha + \beta},
  \label{eq:optimum}
\end{equation}
so $N^* \propto C^{a}$ and $D^* \propto C^{b}$ with $a + b = 1$. The Hoffmann
constants give $a = 0.452$ and $b = 0.548$, parameters and data growing at nearly
the same rate; Kaplan's earlier conclusion that model size should grow much faster
than data came from learning-rate schedules not matched to each run's length, which
Hoffmann corrected. The \term{compute-optimal frontier} is the curve
$(N^*(C), D^*(C))$ in the $(N, D)$ plane, and an \term{IsoFLOP curve} is the loss
along one budget line $6ND = C$, a U in $N$ whose minimum sits on the frontier.
\src{experiments/041-scaling-laws/scripts/scaling\_law\_forms.py, ChinchillaFit.a / .b / .G}

\notesec{The same runs, a different fit}
Besiroglu et al.\ (arXiv:2404.10102)\external{arXiv:2404.10102, not in the
mirror} refit the Approach~3 data and got $E = 1.82$, $A = 482$, $B = 2085$,
$\alpha = 0.35$, $\beta = 0.37$, which moves the allocation to $a = 0.513$ and
$b = 0.487$. The same runs support two frontiers that differ visibly over the
budgets that matter (panel~h), which is the standing argument for reporting an
interval on every exponent rather than a point.

%% SOURCE: experiments/041-scaling-laws/scripts/scaling_law_forms.py (write_constants_table)
%% Published constants typed from the papers named in the Fit records of that script;
%% a and b are derived from alpha and beta there. Never hand-edit this file.
\begin{table}[htbp]
\centering\footnotesize
\caption{Published constants of the two forms, and the allocation exponents they imply. Kaplan's $N_c$ and $D_c$ set the height of a pure power law and have no meaning on their own; Hoffmann's and Besiroglu's rows are two fits of the same Chinchilla training runs, and the difference between them is the point of panel h. Every value is from outside our holdings.}
\label{tab:constants}
\begin{tabular}{llrl}
\toprule
fit & symbol & value & units \\
\midrule
Kaplan 2020 & $N_c$ & $8.8\times 10^{13}$ & parameters \\
 & $\alpha_N$ & $0.076$ & -- \\
 & $D_c$ & $5.4\times 10^{13}$ & tokens \\
 & $\alpha_D$ & $0.095$ & -- \\
\addlinespace
Hoffmann 2022 & $E$ & $1.69$ & loss \\
 & $A$ & $406.4$ & loss $\cdot$ parameters$^{\alpha}$ \\
 & $B$ & $410.7$ & loss $\cdot$ tokens$^{\beta}$ \\
 & $\alpha$ & $0.34$ & -- \\
 & $\beta$ & $0.28$ & -- \\
 & $a = \beta / (\alpha + \beta)$ & $0.452$ & derived \\
 & $b = \alpha / (\alpha + \beta)$ & $0.548$ & derived \\
\addlinespace
Besiroglu 2024 & $E$ & $1.8172$ & loss \\
 & $A$ & $482.01$ & loss $\cdot$ parameters$^{\alpha}$ \\
 & $B$ & $2085.43$ & loss $\cdot$ tokens$^{\beta}$ \\
 & $\alpha$ & $0.3478$ & -- \\
 & $\beta$ & $0.3658$ & -- \\
 & $a = \beta / (\alpha + \beta)$ & $0.513$ & derived \\
 & $b = \alpha / (\alpha + \beta)$ & $0.487$ & derived \\
\bottomrule
\end{tabular}
\end{table}


\notesec{Symbols}
Three names in the fitting code of Section~\ref{sec:ablation} clash with the
symbols above and are listed here so the clash is visible. Table~\ref{tab:symbols}
is hand-authored.

%% SOURCE: hand-authored; the symbols of eqs. 1-4 and of the estimator in
%% experiments/041-scaling-laws/scripts/scaling_law_forms.py (fit_floored).
\begin{table}[htbp]
\centering\footnotesize
\caption{Every symbol in Equations~\ref{eq:kaplan} to \ref{eq:optimum} and in the
estimator, with the three names that collide.}
\label{tab:symbols}
\begin{tabular}{lp{0.78\textwidth}}
\toprule
symbol & meaning \\
\midrule
$L$ & test loss: average loss on a fixed, separate held-out set \\
$N$ & trainable parameters; Kaplan counts non-embedding parameters only \\
$D$ & training data the optimizer consumed, in the unit the loss averages over (tokens, or records here) \\
$C$ & training compute in floating-point operations \\
$6$ & FLOPs per parameter per training example for a transformer, forward plus backward \\
$N_c$, $D_c$ & Kaplan's fitted scale constants, units of parameters and of data \\
$\alpha_N$, $\alpha_D$ & Kaplan's single-axis exponents \\
$E$ & irreducible loss, the floor \\
$A$, $B$ & coefficients of the finite-size and finite-data terms \\
$\alpha$, $\beta$ & exponents of the joint fit \\
$N^*$, $D^*$ & compute-optimal size and data for a budget $C$ \\
$a$, $b$ & allocation exponents, $N^* \propto C^a$, $D^* \propto C^b$, $a + b = 1$ \\
$G$ & prefactor of the optimum, Equation~\ref{eq:optimum} \\
\addlinespace
\texttt{log\_a}, \texttt{log\_e} & in the estimator: $\log A$ and $\log E$, fit in log space so both stay positive; not the allocation exponents $a$, $b$ \\
\texttt{alpha} & in the estimator: the joint-fit $\alpha$ \\
\texttt{r} & the residual, predicted log-loss minus observed log-loss \\
\texttt{starts} & the grid of initial guesses the optimizer is run from, because the floored fit has several local minima \\
\bottomrule
\end{tabular}
\end{table}
